\documentclass{AeroDesign}
\usepackage[utf8]{inputenc}
\usepackage[T1]{fontenc}
\usepackage{geometry}
\usepackage{graphicx}
\usepackage{float}
\usepackage{hyperref}
\usepackage[english,indonesian]{babel}

\title{Optimasi Pemrosesan Template LaTeX Menggunakan Pendekatan Berbasis AI}
\author{Academic Support Agent}
\date{\today}

\begin{document}
\maketitle

\begin{abstract}

\end{abstract}

\chapter{PENDAHULUAN}


\chapter{TINJAUAN PUSTAKA}


% Original content (commented out):
% % ------------------------------------------------------------------------
% ------------------------------------------------------------------------
% relaTeX - Modelo de relatório da equipe AeroDesign ITA, responsável: saulo
%
% Arquivos logo_equipe.pdf , logo_rodape.pdf e logo_SaeBrasil.pdf são auto-explicativos né?
%           [vai na capa]      [no rodapé]      [canto da capa]
%
% Usamos e abusamos de:
% abnTeX2: Modelo de Relatório Técnico/Acadêmico em conformidade com 
% ABNT NBR 10719:2011 Informação e documentação - Relatório técnico e/ou
% científico - Apresentação
%
% ------------------------------------------------------------------------ 
% ------------------------------------------------------------------------

\documentclass[
	% -- opções da classe memoir --
	12pt,				% tamanho da fonte
	oneside,			% para impressão em verso e anverso. Oposto a oneside
	a4paper,			% tamanho do papel. 
	% -- opções da classe abntex2 --
	%chapter=TITLE,		% títulos de capítulos convertidos em letras maiúsculas
	%section=TITLE,		% títulos de seções convertidos em letras maiúsculas
	%subsection=TITLE,	% títulos de subseções convertidos em letras maiúsculas
	%subsubsection=TITLE,% títulos de subsubseções convertidos em letras maiúsculas
	% -- opções do pacote babel --
	english,			% idioma adicional para hifenização
	french,				% idioma adicional para hifenização
	spanish,			% idioma adicional para hifenização
	brazil,				% o último idioma é o principal do documento
	]{AeroDesign}

\newtheorem{defi}{Defini\c{c}\~ao}
% ---
% PACOTES
% ---
\usepackage{float}
\usepackage{amsmath}
\usepackage{caption, subcaption}
\usepackage{multirow} %Para mesclar linhas na tabela
\usepackage{chngpage} %Para estrapolar as margens em caso de figuras/tablas grandes
% ---
% Pacotes fundamentais 
% ---
\usepackage{placeins}
\usepackage{lmodern}			% Usa a fonte Latin Modern
\usepackage[T1]{fontenc}		% Selecao de codigos de fonte.
\usepackage[utf8]{inputenc}		% Codificacao do documento (conversão automática dos acentos)
\usepackage{indentfirst}		% Indenta o primeiro parágrafo de cada seção.
\usepackage{color}				% Controle das cores
\usepackage{graphicx}			% Inclusão de gráficos
\usepackage{microtype} 			% para melhorias de justificação
\usepackage{pdfpages}
\usepackage{pstricks}
\usepackage[a4paper, left=1in, right=0.5in, top=0.5in, bottom=0.5in]{geometry}
\graphicspath{{imgs/}}
%##############
\usepackage{tabu}				% usar ambiente tabu no lugar de tabular
\setlength{\tabulinesep}{1.2mm} % Espaçando itens nas celulas da tabela
%##############

% ---
% Pacotes de citações
% ---
\usepackage[brazilian,hyperpageref]{backref}	 % Paginas com as citações na bibl
\usepackage[alf]{abntex2cite}	% Citações padrão ABNT

% --- 
% CONFIGURAÇÕES DE PACOTES
% --- 

% ---
% Configurações do pacote backref
% Usado sem a opção hyperpageref de backref
\renewcommand{\backrefpagesname}{Citado na(s) página(s):~}
% Texto padrão antes do número das páginas
\renewcommand{\backref}{}
% Define os textos da citação
\renewcommand*{\backrefalt}[4]{
	\ifcase #1 %
		Nenhuma citação no texto.%
	\or
		Citado na página #2.%
	\else
		Citado #1 vezes nas páginas #2.%
	\fi}%
% ---

% -------------------------------------------
% -------------------------------------------
% Informações de dados para CAPA e FOLHA DE ROSTO
% Informar o numero da equipe 
% O título é o nome da equipe!
% Area é o assunto do rela
% Colocar todos os membros como autores, pode colocar \at email se quiser..
% ---
\renewcommand{\nrEquipe}{03}
\titulo{Nome da equipe}
\renewcommand{\area}{Projeto conceitual}
\renewcommand{\capitao}{Nome do capitão}
\autor{
	Integrante 1 \\
	Integrante 2 \\
	Integrante 3 \\
	Você já entendeu \\
}
\local{Brasil}
\data{2016, v-1.0}
\instituicao{Instituição de ensino}
\orientador{Prof. Orientador}
\tipotrabalho{Relatório técnico}


% Configurações de aparência do PDF final

% alterando o aspecto da cor azul
\definecolor{blue}{RGB}{41,5,195}

% ---
% Identificando a categoria da equipe pela numeração
% ---

\global\edef\nrEquipepx{\nrEquipe,px}
\ifdim \nrEquipepx>200px
\definecolor{bgEquipe}{RGB}{0,190,100}
\renewcommand{\clEquipe}{Classe Micro}
\else
\definecolor{bgEquipe}{RGB}{0,128,255}
\renewcommand{\clEquipe}{Classe Advanced}
\fi
\ifdim \nrEquipepx<100px
\definecolor{bgEquipe}{RGB}{255,60,0}
\renewcommand{\clEquipe}{Classe Regular}
\fi 
% informações do PDF
\makeatletter
\hypersetup{
     	%pagebackref=true,
		pdftitle={Relatório \@title}, 
		pdfauthor={\@title},
    	pdfsubject={\imprimirpreambulo},
	    pdfcreator={LaTeX with abnTeX2},
		pdfkeywords={AeroDesign}{SAE Brasil}{abntex}{abntex2}{relatório técnico}, 
		colorlinks=true,       		% false: boxed links; true: colored links
    	linkcolor=red,          	% color of internal links
    	citecolor=blue,        		% color of links to bibliography
    	filecolor=magenta,      		% color of file links
		urlcolor=blue,
		bookmarksdepth=4
}
\makeatother
% --- 

% --- 
% Espaçamentos entre linhas e parágrafos 
% --- 

% O tamanho do parágrafo é dado por:
\setlength{\parindent}{1.2cm}

% Controle do espaçamento entre um parágrafo e outro:
\setlength{\parskip}{0.2cm}  % tente também \onelineskip

% Espaçamento de itens dentro da celula
\makeindex
% ----
% Início do documento
% ----
\begin{document}

% Seleciona o idioma do documento (conforme pacotes do babel)
%\selectlanguage{english}
\selectlanguage{brazil}

% Retira espaço extra obsoleto entre as frases.
\frenchspacing 

% ----------------------------------------------------------
% ELEMENTOS TEXTUAIS
% ----------------------------------------------------------

\textual

\imprimircapa

\pdfbookmark[0]{\contentsname}{toc}
\tableofcontents*
\begin{simbolos}
\input{simbolos.tex}
\end{simbolos}

\input{input}
\input{texto}
\input{output}



\postextual

\bibliography{bibliografia}

\end{document}


\end{document}